\begin{table}[t]
\centering
\caption{Median 90\% interval lengths under the matched gain-loss standard}
\label{tab:app_interval_widths}
\begin{tabular}{lrrrrrr}
\toprule
Asset & Windows & XP & Matched VaR & Matched ES & VaR/XP & ES/XP \\
\midrule
SPX & 1011 & 0.76 & 0.63 & 1.06 & 0.83 & 1.40 \\
EUR & 1082 & 0.30 & 0.24 & 0.45 & 0.77 & 1.48 \\
BTC & 1921 & 3.73 & 3.02 & 5.65 & 0.81 & 1.52 \\
ETH & 1243 & 4.50 & 3.73 & 7.06 & 0.83 & 1.57 \\
BNB & 1243 & 4.56 & 3.27 & 7.65 & 0.72 & 1.68 \\
ADA & 1243 & 4.13 & 2.76 & 6.04 & 0.67 & 1.46 \\
\bottomrule
\end{tabular}
\begin{flushleft}
\footnotesize
Notes: The table fixes \(\tau=0.01\), so that \(\Omega=(1-\tau)/\tau=99\),
and computes the quantile level \(\alpha(\tau)\) satisfying
\(\widehat{VaR}_{\alpha(\tau)}=\widehat{XP}_{\tau}\) in each rolling window.
XP is the median length of the plug-in CLT interval for the conditional expectile.
Matched VaR and matched ES are median two-step FHS interval lengths computed using
the feasible asymptotic variances of Gao and Song (2008), including uncertainty in
the QML volatility forecast, uncertainty in the residual quantile or tail mean, and
their covariance. Medians are reported because
quantile intervals depend on the inverse density at the threshold and may be highly
unstable in windows where this density is close to zero. Ratios above one indicate
wider intervals than XP.
Medians use finite, nonnegative feasible variance estimates.
Matched VaR excludes 19 EUR windows with a negative finite-sample plug-in variance.
\end{flushleft}
\end{table}
